% --- Desenho de grafos (execucao passo a passo) ---
\usepackage{tikz}
\usetikzlibrary{arrows.meta,calc,positioning}

% Cores da caixa "Execucao Passo a Passo" (verde-petroleo, distinta das questoes)
\definecolor{corPassoBorda}{HTML}{2E7D6B}
\definecolor{corPassoBg}{HTML}{F3F8F6}
\definecolor{corPassoTitulo}{HTML}{1E5A4C}
% Cores de estado dos desenhos
\definecolor{corArvore}{HTML}{2C5F8A}     % aresta/vertice ja escolhido
\definecolor{corAtual}{HTML}{E07A1F}      % aresta/vertice examinado neste passo
\definecolor{corRejeitada}{HTML}{C0392B}  % aresta rejeitada / problema
\definecolor{corInativa}{HTML}{AAB3BC}    % aresta ainda nao usada

% Caixa "Execucao Passo a Passo": barra lateral grossa, cantos retos, faixa de titulo
\newtcolorbox{passoapasso}[1][]{%
  breakable,
  enhanced,
  sharp corners,
  colback=corPassoBg,
  colframe=corPassoBorda,
  boxrule=0.5pt,
  leftrule=4pt,
  left=10pt, right=8pt, top=6pt, bottom=6pt,
  fonttitle=\sffamily\bfseries\small,
  coltitle=white,
  colbacktitle=corPassoBorda,
  toptitle=3pt, bottomtitle=3pt,
  title={$\blacktriangleright$\ Execução Passo a Passo --- #1}
}

% Estilos TikZ comuns
\tikzset{
  vtx/.style={circle, draw=corSecundaria, fill=white, line width=0.6pt,
              minimum size=6mm, inner sep=0pt, font=\scriptsize\sffamily\bfseries},
  vtxS/.style={vtx, draw=corArvore, fill=corArvore, text=white},
  vtxAtual/.style={vtx, draw=corAtual, fill=corAtual, text=white},
  vtxRej/.style={vtx, draw=corRejeitada, fill=corRejeitada!15, text=corRejeitada},
  ar/.style={draw=corInativa, line width=0.7pt},
  arSel/.style={draw=corArvore, line width=1.8pt},
  arAtual/.style={draw=corAtual, line width=1.8pt},
  arRej/.style={draw=corRejeitada, line width=1.1pt, dash pattern=on 3pt off 2pt},
  arApagada/.style={draw=corInativa!50, line width=0.5pt, dash pattern=on 1pt off 1.5pt},
  dir/.style={-{Stealth[length=2.2mm,width=1.8mm]}},
  peso/.style={fill=corPassoBg, inner sep=0.8pt, font=\scriptsize\sffamily, text=corSecundaria},
  rot/.style={font=\scriptsize\sffamily\bfseries, text=corAtual, inner sep=1pt}
}

% Painel de um passo: titulo, desenho e explicacao curta (3 por linha)
\newcommand{\passo}[3]{%
  \begin{minipage}[t]{0.325\linewidth}\centering
    {\sffamily\footnotesize\bfseries\color{corPassoTitulo}#1}\par\smallskip
    #2\par\smallskip
    \parbox{\linewidth}{\footnotesize\RaggedRight #3}
  \end{minipage}}
% Painel largo (2 por linha), para grafos maiores
\newcommand{\passoLargo}[3]{%
  \begin{minipage}[t]{0.49\linewidth}\centering
    {\sffamily\footnotesize\bfseries\color{corPassoTitulo}#1}\par\smallskip
    #2\par\smallskip
    \parbox{\linewidth}{\footnotesize\RaggedRight #3}
  \end{minipage}}

% Valor alterado no passo (laranja) e valor definitivo (azul)
\newcommand{\novo}[1]{\textcolor{corAtual}{\textbf{#1}}}
\newcommand{\fixo}[1]{\textcolor{corArvore}{\textbf{#1}}}

% Legenda das cores dos desenhos
\newcommand{\legendaPasso}{%
  \begin{center}\footnotesize\sffamily
  \tikz[baseline=-0.6ex]\draw[arSel](0,0)--(0.5,0); escolhida \qquad
  \tikz[baseline=-0.6ex]\draw[arAtual](0,0)--(0.5,0); examinada no passo \qquad
  \tikz[baseline=-0.6ex]\draw[arRej](0,0)--(0.5,0); rejeitada/problema \qquad
  \tikz[baseline=-0.6ex]\draw[ar](0,0)--(0.5,0); ainda não usada\\[2pt]
  \novo{laranja}: valor alterado no passo \qquad \fixo{azul}: valor definitivo
  \end{center}}

% Etapa em linha: desenho a esquerda, tabela/estado a direita (sem precisar rolar)
% \etapa{titulo}{desenho}{conteudo da direita: tabela + 1-2 linhas}
\newcommand{\etapa}[3]{%
  \par\noindent
  \begin{minipage}{\linewidth}
  {\sffamily\footnotesize\bfseries\color{corPassoTitulo}#1}\par\smallskip
  \begin{minipage}[c]{0.40\linewidth}\centering #2\end{minipage}\hfill
  \begin{minipage}[c]{0.58\linewidth}\footnotesize\RaggedRight #3\end{minipage}
  \end{minipage}\par\medskip
  {\color{corPassoBorda!35}\hrule height 0.4pt}\par\medskip}
% Variante com desenho mais largo (0.55 / 0.43)
\newcommand{\etapaLarga}[3]{%
  \par\noindent
  \begin{minipage}{\linewidth}
  {\sffamily\footnotesize\bfseries\color{corPassoTitulo}#1}\par\smallskip
  \begin{minipage}[c]{0.55\linewidth}\centering #2\end{minipage}\hfill
  \begin{minipage}[c]{0.43\linewidth}\footnotesize\RaggedRight #3\end{minipage}
  \end{minipage}\par\medskip
  {\color{corPassoBorda!35}\hrule height 0.4pt}\par\medskip}

